\chapter{Introducci\'on}
\label{cap:introduccion}

\section{Contexto: gesti\'on veterinaria y digitalizaci\'on incompleta}

La administraci\'on de una cl\'inica veterinaria involucra procesos que, en
ausencia de una herramienta digital centralizada, dependen de registros
manuales o en hojas de c\'alculo dispersas: identificaci\'on de mascotas y
due\~nos, historial de atenciones, control de vacunas y programaci\'on de
citas. La literatura reciente sobre el contexto ecuatoriano confirma que
esta carencia no es an\'ecdota: Cede\~no Ochoa, Catuto Murillo \& Rodas-Silva
documentan expl\'icitamente c\'omo la mayor\'ia de cl\'inicas veterinarias
locales carecen de herramientas tecnol\'ogicas de gesti\'on, lo que provoca
p\'erdida y duplicidad de expedientes por manejo manual~\cite{cedeno2021veterinaria}.
Rodr\'iguez et al. documentan un patr\'on similar en otro contexto
geogr\'afico, motivando un sistema de c\'odigo abierto para el mismo
dominio~\cite{rodriguez2024oscrum}. BIOPET nace de ese mismo diagn\'ostico:
la necesidad de un sistema web que centralice identificaci\'on,
autenticaci\'on, y la gesti\'on de mascotas, citas, consultas y vacunas de
una cl\'inica veterinaria, con evidencia emp\'irica de que efectivamente
funciona y es razonablemente seguro, no solo de que fue programado.

\section{Problema abordado}

El problema que BIOPET aborda tiene dos capas. La primera es funcional:
proveer una plataforma web que permita registrar usuarios con distintos
roles (administrador, veterinario, auxiliar, due\~no), administrar mascotas
asociadas a un due\~no, y gestionar el ciclo de atenci\'on veterinaria
(citas, consultas, vacunas) con control de acceso adecuado a cada rol. La
segunda capa, menos visible pero igual de cr\'itica, es de \textbf{evidencia
emp\'irica}: un sistema acad\'emico puede afirmar que ``funciona'' sin que
esa afirmaci\'on est\'e respaldada por pruebas automatizadas, medici\'on de
rendimiento bajo carga, auditor\'ia de seguridad real o percepci\'on de
usabilidad medida con un instrumento estandarizado. BIOPET trata esa
segunda capa como parte del problema mismo a resolver, no como un anexo
opcional: cada afirmaci\'on cuantitativa de este informe cita su fuente de
evidencia real en el repositorio.

\section{Motivaci\'on}

Tres entregas previas (Entrega~1A, Entrega~1B y Tercera Entrega) generaron
retroalimentaci\'on docente real, registrada \'integramente en
\texttt{docs/observaciones/OBSERVACIONES.md}: 15 observaciones formales,
todas cerradas a la fecha de este informe, cubriendo desde la ausencia de un
repositorio accesible hasta la falta de evidencia Lighthouse o de un DOI de
archivado permanente. Esa bitácora de correcciones es, en s\'i misma, parte
de la motivaci\'on de este documento: la Entrega Final no reescribe el
proyecto desde cero, sino que consolida el trabajo ya evaluado, corrigiendo
expl\'icitamente lo se\~nalado y documentando con honestidad lo que sigue
pendiente.

\section{Preguntas de investigaci\'on}
\label{sec:preguntas-investigacion}

Este informe organiza su evaluaci\'on emp\'irica alrededor de tres
preguntas de investigaci\'on, cada una medible con evidencia real ya
recolectada y respondida expl\'icitamente en el
cap\'itulo~\ref{cap:discusion} (secci\'on~\ref{sec:respuestas-rq}):

\begin{description}[nosep]
  \item[RQ1.] \textquestiondown En qu\'e medida BIOPET satisface los
    requisitos funcionales y no funcionales definidos para la Entrega
    Final, y cu\'ales permanecen expl\'icitamente pendientes?
  \item[RQ2.] \textquestiondown Qu\'e resultados presenta BIOPET en
    cobertura de pruebas, rendimiento, usabilidad, calidad web y seguridad
    bajo las evaluaciones emp\'iricas realmente ejecutadas?
  \item[RQ3.] \textquestiondown En qu\'e medida el proceso de despliegue,
    integraci\'on continua, publicaci\'on de la imagen de contenedor y
    archivado permanente permiten reproducir y verificar el artefacto
    final de forma independiente?
\end{description}

\section{Objetivo general}

Consolidar, evaluar emp\'iricamente y documentar acad\'emicamente el sistema
BIOPET en su versi\'on v1.0.0, integrando evidencia real de calidad
funcional, seguridad, rendimiento, usabilidad, mantenibilidad y
reproducibilidad, enmarcada en el modelo de calidad ISO/IEC~25010 y en una
metodolog\'ia de investigaci\'on de ingenier\'ia expl\'icita, sin inventar
resultados ni ocultar limitaciones.

\subsection*{Objetivos espec\'ificos}
\begin{enumerate}[nosep]
  \item Verificar el cumplimiento de los requisitos funcionales y no
    funcionales declarados en el SRS y documentar la arquitectura del
    sistema (incluidos los m\'odulos de Unidad~IV), distinguiendo
    expl\'icitamente lo verificado de lo pendiente. \emph{(RQ1)}
  \item Reportar la evidencia emp\'irica real de pruebas automatizadas,
    cobertura, rendimiento, usabilidad y calidad web recolectada hasta el
    cierre de esta entrega. \emph{(RQ2)}
  \item Auditar la seguridad del sistema desde tres \'angulos
    complementarios (revisi\'on manual OWASP, an\'alisis din\'amico ZAP,
    an\'alisis est\'atico SpotBugs/Find Security Bugs), incluyendo el
    tratamiento expl\'icito del \'unico hallazgo de severidad alta
    detectado. \emph{(RQ2)}
  \item Documentar el despliegue real en producci\'on, la publicaci\'on de
    la imagen de contenedor en GHCR, el archivado permanente en Zenodo y
    la reproducibilidad de punta a punta del artefacto. \emph{(RQ3)}
  \item Discutir los resultados frente a la literatura relacionada
    disponible y declarar honestamente las amenazas a la validez y las
    limitaciones del estudio, incluyendo condiciones a\'un no verificables
    (por ejemplo, operaci\'on sostenida durante per\'iodos prolongados).
    \emph{(RQ1--RQ3)}
\end{enumerate}

\section{Alcance}

El alcance funcional implementado y verificado cubre: registro y
autenticaci\'on de usuarios por cookies seguras; autorizaci\'on por rol y
por propiedad; CRUD completo de mascotas; y los m\'odulos de citas,
consultas y vacunas incorporados en la Unidad~IV, respaldados por seis
rutinas almacenadas de PostgreSQL. Quedan expl\'icitamente \textbf{fuera}
del alcance implementado, y as\'i se documentan en el cap\'itulo de
requisitos (cap\'itulo~\ref{cap:requisitos}): historial cl\'inico
cronol\'ogico completo como m\'odulo independiente, integraci\'on con
dispositivos IoT, facturaci\'on digital, reportes exportables en PDF/Excel,
y recomendaciones cl\'inicas generadas por IA externa. El alcance de
evaluaci\'on emp\'irica cubre pruebas automatizadas, cobertura, rendimiento,
usabilidad, calidad web, y seguridad est\'atica/din\'amica, todas ejecutadas
en un entorno acad\'emico (desarrollo local y despliegue en el plan gratuito
de Render), no en un entorno industrial de alta disponibilidad.

\section{Contribuci\'on del proyecto}

La contribuci\'on principal de BIOPET, frente al panorama de trabajos
relacionados revisado en el cap\'itulo~\ref{cap:trabajos-relacionados}, no
es una t\'ecnica nueva de ingenier\'ia de software, sino la
\textbf{amplitud verificable de la evidencia emp\'irica} recolectada sobre
un mismo artefacto real: cobertura de pruebas, rendimiento bajo carga,
usabilidad medida con instrumento estandarizado, calidad web automatizada, y
seguridad auditada por tres m\'etodos distintos, todo trazable a
requisitos formales y reproducible con comandos documentados
(\texttt{make up}, \texttt{mvn clean verify}, scripts de
\texttt{scripts/}). Ning\'un trabajo relacionado revisado reporta las seis
dimensiones a la vez sobre el mismo sistema (secci\'on~\ref{sec:sintesis-trabajos-relacionados}).

\section{Organizaci\'on de este informe}

El resto del documento se organiza as\'i: el cap\'itulo~\ref{cap:marco-teorico}
resume el marco te\'orico y normativo usado en todo el informe; el
cap\'itulo~\ref{cap:trabajos-relacionados}
revisa los trabajos relacionados bajo un protocolo inspirado en
PRISMA~2020; el cap\'itulo~\ref{cap:metodologia} describe la metodolog\'ia
de investigaci\'on de ingenier\'ia adoptada; el cap\'itulo~\ref{cap:requisitos}
documenta la ingenier\'ia de requisitos; el cap\'itulo~\ref{cap:arquitectura}
la arquitectura y las decisiones de dise\~no; el cap\'itulo~\ref{cap:implementacion}
los m\'odulos implementados; los cap\'itulos~\ref{cap:pruebas-calidad} y
\ref{cap:seguridad} la evidencia de pruebas/calidad y de seguridad; el
cap\'itulo~\ref{cap:despliegue} el despliegue y la reproducibilidad; el
cap\'itulo~\ref{cap:trazabilidad-final} la matriz de trazabilidad; los
cap\'itulos~\ref{cap:resultados-final} y \ref{cap:discusion} integran y
discuten los resultados; el cap\'itulo~\ref{cap:amenazas-limitaciones-final}
documenta las amenazas a la validez y las limitaciones; el
cap\'itulo~\ref{cap:conclusiones-final} presenta las conclusiones; el
cap\'itulo~\ref{cap:declaraciones} recoge las declaraciones formales
(CRediT, \'etica, disponibilidad); y los anexos
(cap\'itulo~\ref{cap:anexos-final}) complementan con glosario, endpoints y
documentos de referencia.
